%% ch-forth-programming.tex — FORTH Programming Patterns and Style
%% SOURCE: .claude/CLAUDE.md, src/word_source/
%% Included from user-guide/main.tex

%%━━━━━━━━━━━━━━━━━━━━━━━━━━━━━━━━━━━━━━━━━━━━━━━━━━━━━━━━━━━━━━━━━━━━━━━━━
\part{FORTH Programming}
%%━━━━━━━━━━━━━━━━━━━━━━━━━━━━━━━━━━━━━━━━━━━━━━━━━━━━━━━━━━━━━━━━━━━━━━━━━

%%────────────────────────────────────────────────────────────────────────────
\chapter{Thinking in FORTH}
\label{chap:thinking}
%%────────────────────────────────────────────────────────────────────────────

FORTH is a different kind of language. If you come from C, Python, or Java, a
few shifts in thinking will make it click faster.

%% ─────────────────────────────────────────────────────────────────────────
\section{Postfix Notation}
\label{sec:postfix}

Every language you have probably used puts operators between operands:
\texttt{3 + 4}. FORTH puts the operands first and the operator last:
\texttt{3 4 +}. This is postfix (or reverse Polish) notation.

The reason is the stack. When you type \texttt{3 4 +}, the runtime pushes 3,
pushes 4, then \texttt{+} pops both and pushes 7. There are no parentheses
for grouping because order of operations is completely explicit. Compare:

\begin{center}
\begin{tabular}{ll}
\toprule
Infix (C) & FORTH \\
\midrule
\texttt{(3 + 4) * 2} & \texttt{3 4 + 2 *} \\
\texttt{a * a + b * b} & \texttt{a a * b b * +} \\
\texttt{(x - 1) / (x + 1)} & \texttt{x 1 - x 1 + /} \\
\bottomrule
\end{tabular}
\end{center}

You evaluate left to right, and the result of each operation is immediately
available for the next.

%% ─────────────────────────────────────────────────────────────────────────
\section{Words Are Verbs}
\label{sec:words-are-verbs}

In FORTH, a ``word'' is any named operation --- primitive or user-defined.
Numbers push themselves; everything else is a word that does something to the
stack. \texttt{DUP} is a word. \texttt{+} is a word. \texttt{HELLO} (once you
define it) is a word.

This uniformity means the language is extensible in a deep way. After you
define \texttt{SQUARE}, it is indistinguishable from a primitive --- it appears
in \texttt{WORDS}, it runs when you type it, and other words can call it. The
line between ``built-in'' and ``user-defined'' is stylistic, not structural.

%% ─────────────────────────────────────────────────────────────────────────
\section{Factor Early, Factor Often}
\label{sec:factor}

The FORTH style is to decompose programs into many small words, each doing one
thing. A word with more than five or six lines in its body is usually a sign
that it wants to be split.

\begin{lstlisting}[language=Forth]
( Bad: one big word )
: PROCESS-PAYMENT  ( amount account-id -- )
    DUP VALID-ACCOUNT? IF
        OVER CHECK-FUNDS IF
            SWAP DEDUCT
            LOG-TRANSACTION
        ELSE ." Insufficient funds" CR THEN
    ELSE ." Invalid account" CR THEN
    DROP ;
\end{lstlisting}

\begin{lstlisting}[language=Forth]
( Better: factored )
: REJECT-INVALID  ." Invalid account" CR DROP DROP ;
: REJECT-FUNDS    ." Insufficient funds" CR DROP DROP ;
: COMMIT-PAYMENT  ( amount account-id -- )
    SWAP DEDUCT LOG-TRANSACTION ;

: PROCESS-PAYMENT  ( amount account-id -- )
    DUP VALID-ACCOUNT? NOT IF  REJECT-INVALID  EXIT  THEN
    OVER CHECK-FUNDS  NOT IF  REJECT-FUNDS     EXIT  THEN
    COMMIT-PAYMENT ;
\end{lstlisting}

Factored code is easier to test interactively, easier to reuse, and easier to
read. You can test \texttt{COMMIT-PAYMENT} in isolation before wiring it into
\texttt{PROCESS-PAYMENT}.

%% ─────────────────────────────────────────────────────────────────────────
\section{Stack Comments Are Documentation}
\label{sec:stack-comments}

FORTH comments use parentheses: \texttt{( this is a comment )}. By
convention, every word definition begins with a stack-effect comment:

\begin{lstlisting}[language=Forth]
: AREA  ( width height -- area )
    * ;
\end{lstlisting}

The stack comment is the primary documentation. Write it first, before the
body, and treat it as a contract: the word must consume exactly the items
before the \texttt{--} and produce exactly the items after. Violations cause
hard-to-debug stack confusion.

StarForth does not enforce stack effects at runtime --- it is your
responsibility to make them accurate.

%% ─────────────────────────────────────────────────────────────────────────
\section{The Interactive Development Loop}
\label{sec:dev-loop}

FORTH's interactive nature is a feature, not a limitation. The typical
development cycle is:

\begin{enumerate}
  \item Type a small fragment at the \texttt{ok>} prompt and test it.
  \item When it works, wrap it in a \texttt{: NAME ... ;} definition.
  \item Test the word interactively.
  \item Build larger words on top of working smaller ones.
\end{enumerate}

For example, building a Fahrenheit-to-Celsius converter step by step:

\begin{lstlisting}[language=Forth]
ok> 32 .
32 ok>

ok> 212 32 - .
180 ok>

ok> 212 32 - 5 * 9 / .
100 ok>

ok> : F>C  ( f -- c )  32 - 5 * 9 / ;
ok> 212 F>C .
100 ok>
ok> 32 F>C .
0 ok>
ok> 98 F>C .
36 ok>
\end{lstlisting}

Each step confirmed before moving on.

%%────────────────────────────────────────────────────────────────────────────
\chapter{Common Programming Patterns}
\label{chap:patterns}
%%────────────────────────────────────────────────────────────────────────────

%% ─────────────────────────────────────────────────────────────────────────
\section{Accumulator Pattern}
\label{sec:pattern-accumulator}

Use a variable as an accumulator and update it in a loop:

\begin{lstlisting}[language=Forth]
VARIABLE SUM

: SUM-TO  ( n -- sum )
    0 SUM !
    1 + 1 DO
        I SUM @ + SUM !
    LOOP
    SUM @ ;
\end{lstlisting}

\begin{lstlisting}[language=Forth]
ok> 10 SUM-TO .
55 ok>
\end{lstlisting}

The pattern: initialize the variable before the loop, update inside, read
after.

%% ─────────────────────────────────────────────────────────────────────────
\section{Table Dispatch}
\label{sec:pattern-dispatch}

Use \texttt{EXECUTE} to implement a jump table. First, build an array of
execution tokens:

\begin{lstlisting}[language=Forth]
: GREET-EN  ." Hello" ;
: GREET-ES  ." Hola" ;
: GREET-FR  ." Bonjour" ;

CREATE GREET-TABLE
    ' GREET-EN ,
    ' GREET-ES ,
    ' GREET-FR ,

: GREET  ( lang-index -- )
    CELLS GREET-TABLE + @ EXECUTE  CR ;
\end{lstlisting}

\begin{lstlisting}[language=Forth]
ok> 0 GREET
Hello ok>
ok> 1 GREET
Hola ok>
ok> 2 GREET
Bonjour ok>
\end{lstlisting}

\texttt{CELLS} converts the index to a byte offset. \texttt{EXECUTE} runs the
stored execution token.

%% ─────────────────────────────────────────────────────────────────────────
\section{Counted String Pattern}
\label{sec:pattern-counted-string}

A common FORTH idiom is the counted string: a byte at the start holds the
length, followed by the characters. \texttt{COUNT} converts a counted-string
address to the \texttt{( addr len )} form that \texttt{TYPE} expects:

\begin{lstlisting}[language=Forth]
: PRINT-COUNTED  ( c-addr -- )
    COUNT TYPE CR ;
\end{lstlisting}

\texttt{S"} produces an \texttt{( addr len )} pair directly:

\begin{lstlisting}[language=Forth]
ok> S" hello world" TYPE CR
hello world
ok>
\end{lstlisting}

%% ─────────────────────────────────────────────────────────────────────────
\section{Guard Clauses}
\label{sec:pattern-guard}

Use \texttt{EXIT} at the start of a word to reject bad inputs early. This is
cleaner than deeply nested \texttt{IF}s:

\begin{lstlisting}[language=Forth]
: SAFE-DIVIDE  ( a b -- a/b )
    DUP 0= IF
        ." Division by zero" CR
        DROP DROP  0
        EXIT
    THEN
    / ;
\end{lstlisting}

\begin{lstlisting}[language=Forth]
ok> 10 2 SAFE-DIVIDE .
5 ok>
ok> 10 0 SAFE-DIVIDE .
Division by zero
0 ok>
\end{lstlisting}

%% ─────────────────────────────────────────────────────────────────────────
\section{State Machine in FORTH}
\label{sec:pattern-state-machine}

Encode state as an integer in a variable, and dispatch on it:

\begin{lstlisting}[language=Forth]
0 CONSTANT STATE-IDLE
1 CONSTANT STATE-RUN
2 CONSTANT STATE-DONE

VARIABLE MACHINE-STATE
STATE-IDLE MACHINE-STATE !

: TRANSITION  ( new-state -- )  MACHINE-STATE ! ;
: STATE@      ( -- state )      MACHINE-STATE @ ;

: STEP
    STATE@ CASE
        STATE-IDLE OF  ." Idle -> Run"  STATE-RUN  TRANSITION  ENDOF
        STATE-RUN  OF  ." Running..."   STATE-DONE TRANSITION  ENDOF
        STATE-DONE OF  ." Done."                                ENDOF
    ENDCASE
    CR ;
\end{lstlisting}

\texttt{CASE} / \texttt{OF} / \texttt{ENDOF} / \texttt{ENDCASE} are the
FORTH multi-way branch words.

%%────────────────────────────────────────────────────────────────────────────
\chapter{Debugging at the Prompt}
\label{chap:debugging}
%%────────────────────────────────────────────────────────────────────────────

%% ─────────────────────────────────────────────────────────────────────────
\section{Using .S Liberally}
\label{sec:debug-dots}

\texttt{.S} is non-destructive --- it leaves the stack intact. Sprinkle it
freely while working out what a piece of code does:

\begin{lstlisting}[language=Forth]
ok> 5 7 .S * .S .
<2> 5 7
<1> 35
35 ok>
\end{lstlisting}

You can also insert \texttt{.S} inside a word definition during debugging (it
compiles to a call to the runtime print-stack routine):

\begin{lstlisting}[language=Forth]
: BUGGY  ( a b -- result )
    .S        \ what do we have?
    +
    .S        \ what did + produce?
    2 * ;
\end{lstlisting}

Remove the \texttt{.S} calls once the word works correctly.

%% ─────────────────────────────────────────────────────────────────────────
\section{SEE: Decompiling a Word}
\label{sec:debug-see}

\texttt{SEE} decompiles a word, showing its compiled form:

\begin{lstlisting}[language=Forth]
ok> SEE DOUBLE
: DOUBLE  DUP + ; ok>
\end{lstlisting}

The output is a reconstruction from the compiled threaded code. It may not
look identical to the original source but shows the same operations.

%% ─────────────────────────────────────────────────────────────────────────
\section{WORDS: Finding What Exists}
\label{sec:debug-words}

When you cannot remember whether you defined a word:

\begin{lstlisting}[language=Forth]
ok> WORDS
\end{lstlisting}

The output is long. You can pipe StarForth output through \texttt{grep} at
the shell level:

\begin{lstlisting}[language=bash]
echo "WORDS BYE" | ./build/amd64/standard/starforth | grep SQUARE
\end{lstlisting}

%% ─────────────────────────────────────────────────────────────────────────
\section{Common Mistakes}
\label{sec:debug-mistakes}

\paragraph{Wrong argument order.}
Postfix is different from most calculators. \texttt{10 3 -} gives 7, not
\(-7\). When in doubt, use \texttt{.S} before and after each operation.

\paragraph{Forgetting to consume values.}
If \texttt{.S} shows values you did not expect, a word left extra values on
the stack. Check every word's documented stack effect.

\paragraph{Control flow outside a definition.}
\texttt{IF}, \texttt{DO}, \texttt{BEGIN}, etc.\ are compile-only. Typing them
at the prompt produces an error. Wrap them inside \texttt{: ... ;}.

\paragraph{Mismatched IF / THEN.}
Every \texttt{IF} needs a \texttt{THEN}. Forgetting one leaves the compiler in
a confused state. An error mid-definition usually means the definition was not
completed; start over with a fresh \texttt{:} definition.

\paragraph{Return stack imbalance.}
Every \texttt{>R} in a word must be matched by an \texttt{R>} before
\texttt{;}. Leaving items on the return stack corrupts the call chain.

%%────────────────────────────────────────────────────────────────────────────
\chapter{Working with Numbers}
\label{chap:numbers-detail}
%%────────────────────────────────────────────────────────────────────────────

%% ─────────────────────────────────────────────────────────────────────────
\section{Integer Ranges}
\label{sec:int-ranges}

A StarForth cell is a 64-bit signed integer on a 64-bit host. The range is
$-2^{63}$ to $2^{63}-1$ (approximately $\pm9.2 \times 10^{18}$). Overflow
wraps silently (two's complement arithmetic) --- StarForth does not trap integer
overflow by default.

\begin{lstlisting}[language=Forth]
ok> 9223372036854775807 1 + .
-9223372036854775808 ok>
\end{lstlisting}

That is the maximum signed 64-bit integer wrapping to the minimum. Keep this
in mind when working with large values.

%% ─────────────────────────────────────────────────────────────────────────
\section{Numeric Literals}
\label{sec:numeric-literals}

StarForth accepts numbers in the current base. In decimal (the default):

\begin{lstlisting}[language=Forth]
ok> 255 .
255 ok>
\end{lstlisting}

After \texttt{HEX}:

\begin{lstlisting}[language=Forth]
ok> HEX
ok> FF .
255 ok>
ok> DECIMAL
\end{lstlisting}

Negative numbers use the minus sign prefix:

\begin{lstlisting}[language=Forth]
ok> -7 ABS .
7 ok>
\end{lstlisting}

%% ─────────────────────────────────────────────────────────────────────────
\section{Double-Cell Numbers}
\label{sec:double-cell}

When you need values larger than one cell, FORTH uses double-cell numbers ---
two consecutive cells treated as a single 128-bit value (on a 64-bit host).
Words that produce or consume double-cells are conventionally prefixed with
\texttt{D} (e.g.\ \texttt{D+}, \texttt{D-}) or suffixed with \texttt{D} in
their stack comments. A double cell occupies two stack positions: the lower
(less significant) cell on top.

\texttt{M*} multiplies two single cells and produces a double cell result:

\begin{lstlisting}[language=Forth]
ok> 1000000 1000000 M* D. CR
1000000000000
ok>
\end{lstlisting}

Double-cell arithmetic is used when single-cell overflow is possible.

%% ─────────────────────────────────────────────────────────────────────────
\section{Pictured Numeric Output}
\label{sec:pictured-output}

FORTH provides a formatting mechanism called pictured numeric output for
constructing number strings character by character:

\begin{lstlisting}[language=Forth]
: HEX-PRINT  ( n -- )
    BASE @ >R  HEX
    <# # # # # # # # # #> TYPE CR
    R> BASE ! ;
\end{lstlisting}

The sequence \texttt{<\#} ... \texttt{\#>} builds the number string in a
scratch buffer. \texttt{\#} converts one digit; \texttt{\#S} converts all
remaining digits; \texttt{HOLD} inserts a literal character; \texttt{\#>}
finishes and returns the string address and length.

A practical example --- printing a number with a comma separator:

\begin{lstlisting}[language=Forth]
: .COMMA  ( n -- )
    <#
    #S
    HOLD 44    \ 44 = comma character
    #S
    #>
    TYPE CR ;
\end{lstlisting}

%%────────────────────────────────────────────────────────────────────────────
\chapter{Vocabularies and the Dictionary}
\label{chap:vocabularies}
%%────────────────────────────────────────────────────────────────────────────

%% ─────────────────────────────────────────────────────────────────────────
\section{How the Dictionary Is Organized}
\label{sec:dict-structure}

The StarForth dictionary is a linked list of word entries. The newest word is
at the head; \texttt{FIND} searches from head toward tail. This means a
newly-defined word with the same name as an existing word shadows the older
one. The older word is not deleted --- it is simply unreachable through the
normal search path until the newer definition is removed with \texttt{FORGET}.

\begin{lstlisting}[language=Forth]
ok> : HELLO  ." Version 1" CR ;
ok> HELLO
Version 1
ok> : HELLO  ." Version 2" CR ;
ok> HELLO
Version 2
ok> FORGET HELLO
ok> HELLO
Version 1
ok>
\end{lstlisting}

\texttt{FORGET} removes the named word and all words defined after it. Use it
carefully --- definitions made after the word you want to forget will also
disappear.

%% ─────────────────────────────────────────────────────────────────────────
\section{VOCABULARY}
\label{sec:vocabulary}

A vocabulary is a named wordlist. StarForth supports multiple vocabularies,
allowing you to organize related words and avoid name collisions.

\begin{lstlisting}[language=Forth]
VOCABULARY GRAPHICS
GRAPHICS DEFINITIONS

: DRAW-LINE   ... ;
: DRAW-CIRCLE ... ;

FORTH DEFINITIONS
\end{lstlisting}

After \texttt{GRAPHICS DEFINITIONS}, new words go into the \texttt{GRAPHICS}
vocabulary. \texttt{FORTH DEFINITIONS} returns to defining into the FORTH
vocabulary.

To make a vocabulary searchable, name it in the search order. The
\texttt{ALSO} word adds a vocabulary to the front of the search order:

\begin{lstlisting}[language=Forth]
ok> ALSO GRAPHICS
ok> DRAW-LINE   \ now visible
\end{lstlisting}

%% ─────────────────────────────────────────────────────────────────────────
\section{SEARCH-WORDLIST}
\label{sec:search-wordlist}

\texttt{SEARCH-WORDLIST} searches a specific wordlist by name and count,
returning an execution token and flags if found, or zero if not found:

\begin{lstlisting}[language=Forth]
ok> S" DRAW-LINE" FORTH-WORDLIST SEARCH-WORDLIST
\end{lstlisting}

This is useful when you want to test for word existence programmatically
before attempting to execute it.

%%────────────────────────────────────────────────────────────────────────────
\chapter{Working with Text and I/O}
\label{chap:text-io}
%%────────────────────────────────────────────────────────────────────────────

%% ─────────────────────────────────────────────────────────────────────────
\section{Reading User Input}
\label{sec:user-input}

\texttt{ACCEPT} reads a line from the terminal:

\begin{lstlisting}[language=Forth]
CREATE INPUT-BUF 80 ALLOT

: GET-LINE  ( -- addr count )
    INPUT-BUF 80 ACCEPT ;

: ECHO-LINE
    ." Enter text: "
    GET-LINE TYPE CR ;
\end{lstlisting}

\begin{lstlisting}[language=Forth]
ok> ECHO-LINE
Enter text: Hello, FORTH!
Hello, FORTH!
ok>
\end{lstlisting}

\texttt{KEY} reads a single character without echoing:

\begin{lstlisting}[language=Forth]
: WAIT-FOR-ENTER
    ." Press Enter to continue..."
    BEGIN KEY 13 = UNTIL CR ;
\end{lstlisting}

%% ─────────────────────────────────────────────────────────────────────────
\section{Parsing Input}
\label{sec:parsing}

\texttt{PARSE} scans the input stream up to a delimiter and returns address
and count. \texttt{WORD} scans to the next whitespace-delimited token and
stores it as a counted string:

\begin{lstlisting}[language=Forth]
\ Parse everything up to the end of line
: PARSE-LINE  ( -- addr count )
    10 PARSE ;   \ 10 = newline
\end{lstlisting}

\texttt{WORD} is used internally by the interpreter to tokenize input. At
the prompt, you rarely need to call it directly.

%% ─────────────────────────────────────────────────────────────────────────
\section{Number Conversion}
\label{sec:number-convert}

\texttt{NUMBER} (or \texttt{NUMBER?} in some implementations) attempts to
parse a string as a number. In StarForth, the outer interpreter does this
automatically when a word is not found in the dictionary. For programmatic
parsing, use the base-aware conversion words.

%% ─────────────────────────────────────────────────────────────────────────
\section{Formatted Output}
\label{sec:formatted-output}

Build formatted output using the print words together:

\begin{lstlisting}[language=Forth]
: PRINT-ROW  ( value label-addr label-len -- )
    TYPE           \ print label
    ." : "
    .R  ;          \ print value right-justified (needs width on stack)

: SHOW-STATS
    ." --- Statistics ---" CR
    42  S" Iterations" 7 ROLL .R  CR
    100 S" Percent"    7 ROLL .R  CR ;
\end{lstlisting}

For tables, use \texttt{.R} (right-justified print) to align columns:

\begin{lstlisting}[language=Forth]
: PRINT-TABLE
    ." N     N^2   N^3" CR
    10 1 DO
        I     5 .R
        I I * 7 .R
        I I * I * 9 .R
        CR
    LOOP ;
\end{lstlisting}

\begin{lstlisting}[language=Forth]
ok> PRINT-TABLE
    N     N^2      N^3
    1       1        1
    2       4        8
    3       9       27
    4      16       64
    5      25      125
    6      36      216
    7      49      343
    8      64      512
    9      81      729
ok>
\end{lstlisting}

%%────────────────────────────────────────────────────────────────────────────
\chapter{Writing and Loading Source Files}
\label{chap:source-files}
%%────────────────────────────────────────────────────────────────────────────

%% ─────────────────────────────────────────────────────────────────────────
\section{The .4th File Convention}
\label{sec:4th-files}

StarForth source files use the \texttt{.4th} extension by convention. A
\texttt{.4th} file is plain text containing FORTH words and definitions,
exactly as you would type them at the prompt:

\begin{lstlisting}[language=Forth]
\ mylib.4th
\ A small library of utility words

: MAX3  ( a b c -- max )
    ROT MAX MAX ;

: MIN3  ( a b c -- min )
    ROT MIN MIN ;

: CLAMP  ( n lo hi -- n' )
    ROT        \ ( lo hi n )
    MIN        \ ( lo n )
    MAX ;      \ ( result )
\end{lstlisting}

Load it with \texttt{INCLUDE}:

\begin{lstlisting}[language=Forth]
ok> S" mylib.4th" INCLUDE
ok> 5 1 10 CLAMP .
5 ok>
ok> 0 1 10 CLAMP .
1 ok>
ok> 99 1 10 CLAMP .
10 ok>
\end{lstlisting}

%% ─────────────────────────────────────────────────────────────────────────
\section{Comments in Source Files}
\label{sec:comments}

FORTH supports two comment forms:

\begin{lstlisting}[language=Forth]
\ This is a line comment (backslash-space)
( This is a paren comment )
\end{lstlisting}

The backslash comment extends to end of line. The paren comment extends to the
closing parenthesis. Both are legal anywhere --- in the REPL, in definitions,
or in source files. By convention, stack effects use paren comments and
explanations use backslash comments:

\begin{lstlisting}[language=Forth]
: HYPOTENUSE  ( a b -- c )
    \ computes sqrt(a^2 + b^2) using integer approximation
    DUP *  SWAP DUP *  + ;  ( a^2+b^2 -- needs SQRT )
\end{lstlisting}

%% ─────────────────────────────────────────────────────────────────────────
\section{Conditional Compilation}
\label{sec:conditional-compile}

Use \texttt{[IF]} / \texttt{[ELSE]} / \texttt{[THEN]} for compile-time
conditionals. These are immediate words that run at compile time rather than
being compiled:

\begin{lstlisting}[language=Forth]
\ Define a debug-print word only in debug builds
HAS-DEBUG [IF]
    : DPRINT  ( n -- )  ." DEBUG: " . CR ;
[ELSE]
    : DPRINT  ( n -- )  DROP ;   \ no-op in production
[THEN]
\end{lstlisting}

%% ─────────────────────────────────────────────────────────────────────────
\section{Loading Capsule Files}
\label{sec:loading-capsules}

StarForth extends the standard FORTH file-loading mechanism with the
\texttt{EXEC} word, which loads named capsule files by name from the capsule
directory:

\begin{lstlisting}[language=Forth]
S" ACL.4th" EXEC      \ load the ACL subsystem
S" doe.4th" EXEC      \ load the DoE experiment harness
\end{lstlisting}

Capsules are content-addressed: the capsule ID is its XXHash64 content hash.
Any modification to the capsule changes its hash, and the system detects the
change at load time.

%%────────────────────────────────────────────────────────────────────────────
\chapter{The Test Suite}
\label{chap:tests}
%%────────────────────────────────────────────────────────────────────────────

%% ─────────────────────────────────────────────────────────────────────────
\section{Running the Tests}
\label{sec:running-tests}

StarForth ships 936+ tests covering every word category:

\begin{lstlisting}[language=bash]
make test
\end{lstlisting}

Tests are organized in POST (Power-On Self Test) order:

\begin{enumerate}
  \item Unit tests: Q48.16 fixed-point, inference statistics, decay slope
    inference.
  \item Dictionary tests: FORTH-79 word validation across 22 categories.
  \item Integration tests and stress tests.
  \item Adversarial and fuzzing tests (\texttt{break\_me\_tests.c}).
\end{enumerate}

Additional options:

\begin{lstlisting}[language=bash]
make test FAIL_FAST=1    # Stop on first failure
make test PROFILE=1      # With basic profiling
\end{lstlisting}

%% ─────────────────────────────────────────────────────────────────────────
\section{Writing Your Own Tests}
\label{sec:writing-tests}

A minimal test pattern in FORTH is:

\begin{lstlisting}[language=Forth]
: TEST-EQUAL  ( actual expected -- )
    = IF ." PASS" ELSE ." FAIL" THEN CR ;

\ Test SQUARE
9 SQUARE 81 TEST-EQUAL
3 SQUARE  9 TEST-EQUAL
0 SQUARE  0 TEST-EQUAL
\end{lstlisting}

For larger test suites, use a counter:

\begin{lstlisting}[language=Forth]
VARIABLE TEST-PASS
VARIABLE TEST-FAIL
0 TEST-PASS !
0 TEST-FAIL !

: CHECK  ( actual expected "desc" -- )
    PARSE-LINE  TYPE ." : "
    = IF
        ." PASS"
        TEST-PASS @ 1 + TEST-PASS !
    ELSE
        ." FAIL"
        TEST-FAIL @ 1 + TEST-FAIL !
    THEN  CR ;

: REPORT
    ." Passed: " TEST-PASS @ . CR
    ." Failed: " TEST-FAIL @ . CR ;

\ Usage:
9 SQUARE 81 CHECK square of 9
0 SQUARE  0 CHECK square of 0
REPORT
\end{lstlisting}
